% ============================================================================
%  2. CÁC CÔNG TRÌNH LIÊN QUAN
% ============================================================================
\section{Các công trình liên quan}
\label{sec:lien-quan}

\subsection{Khảo sát các hướng tiếp cận hiện tại}

Các công trình về hỏi đáp và truy hồi văn bản pháp luật có thể xếp vào bốn hướng, khác
nhau ở chỗ \emph{tri thức pháp lý nằm ở đâu}.

\subsubsection{Hướng 1 --- Truy hồi ngữ nghĩa bằng mạng nơ-ron}

Tri thức nằm trong tham số của mô hình ngôn ngữ. PhoBERT~\cite{phobert} là xương sống cho
phần lớn hệ truy hồi luật tiếng Việt; Pham và cộng sự~\cite{phamMultistage} đề xuất truy
hồi nhiều giai đoạn cho văn bản luật tiếng Việt; Tiên và cộng sự~\cite{tienSynthetic} dùng
dữ liệu tổng hợp do mô hình lớn sinh ra để khắc phục khan hiếm nhãn. Ở quy mô quốc tế,
LEGAL-BERT~\cite{legalbert} cho thấy tiếp tục tiền huấn luyện theo miền pháp luật cải thiện
rõ các tác vụ pháp lý, còn Nguyen và cộng sự~\cite{nguyenAttentive} đưa kiến trúc chú ý
phân cấp cho văn bản luật dài. Điểm chung: đầu ra là một \emph{đoạn văn bản} được xếp hạng,
không phải tri thức có cấu trúc, nên hệ thống không giải thích được vì sao điều khoản đó áp
dụng và không tính được mức phạt.

\subsubsection{Hướng 2 --- Ontology và đồ thị tri thức pháp luật}

Tri thức được mô hình hoá tường minh. Legal-Onto~\cite{legalonto} đề xuất mô hình tri thức
gồm khái niệm có cấu trúc, ba loại quan hệ, luật suy diễn và đồ thị cụm từ khoá, áp dụng
cho Luật Đất đai; đây chính là dòng nghiên cứu mà mô hình $K$ của đồ án kế thừa, cùng với
nền tảng Rela-model~\cite{relamodel} và COKB~\cite{cokb}. Các mở rộng sau đó phủ thêm luật
lao động~\cite{phamLabor,phamLegalOntoUpdate}, tích hợp ontology với đồ thị tri thức cho
truy vấn pháp lý~\cite{phamOntoKG}, và gần đây là kiến trúc ontology ba tầng đa
miền~\cite{phamUnified}. Đáng chú ý nhất với đồ án là Dang và cộng sự~\cite{dangKG} ---
công trình đồ thị tri thức cho chính luật giao thông đường bộ Việt Nam, nhưng dựa trên Luật
23/2008/QH12 đã hết hiệu lực.

\subsubsection{Hướng 3 --- RAG và mô hình ngôn ngữ lớn cho miền pháp luật}

Đây là hướng phổ biến nhất hiện nay, và cũng là hướng nhóm \emph{chủ động không chọn}. Dahl
và cộng sự~\cite{dahl} đo được LLM đa dụng bịa căn cứ pháp lý một cách có hệ thống; Magesh
và cộng sự~\cite{magesh} đánh giá tiền đăng ký các công cụ RAG pháp lý thương mại và ghi
nhận tỉ lệ ảo giác 17--33\% dù nhà cung cấp quảng cáo ``hallucination-free''; benchmark
ALCE~\cite{alce} cho thấy trích dẫn do mô hình sinh thường không trung thực với nguồn. Ngay
cả các hệ kết hợp đồ thị tri thức với RAG cho luật Việt Nam~\cite{phamKGRAG,laura} vẫn giữ
mô hình sinh ở khâu cuối, nên vẫn thừa hưởng rủi ro này.

\subsubsection{Hướng 4 --- Hiệu lực theo thời gian của quy phạm}

Lý thuyết đã tương đối đầy đủ ở quốc tế: Governatori và cộng sự~\cite{governatori} hình
thức hoá ba loại sửa đổi (thay thế, bãi bỏ một phần, huỷ bỏ hồi tố) bằng logic khả bác có
gắn thời gian; Palmirani và cộng sự~\cite{palmirani} quản lý tiến hoá văn bản theo ba chiều
hiệu lực; chuẩn LegalRuleML~\cite{legalruleml} đưa toán tử nghĩa vụ cùng chiều thời gian vào
biểu diễn; de Martim~\cite{demartim} tách tác phẩm trừu tượng khỏi biểu hiện có phiên bản để
trả lời tất định câu hỏi ``văn bản tại thời điểm $t$''. Phía Việt Nam, Pham và cộng
sự~\cite{phamLegalOntoUpdate} mới dừng ở đối sánh hai phiên bản Luật Đất đai 2013 và 2024
--- tức phát hiện thay đổi, chưa phải truy vấn theo mốc thời gian.

\subsection{Các bộ dữ liệu đã có}

\begin{table}[H]
\centering
\caption{Các bộ dữ liệu hỏi đáp và truy hồi pháp luật tiêu biểu}
\label{tab:datasets}
\small
\renewcommand{\arraystretch}{1.25}
\begin{tabularx}{\textwidth}{@{}>{\raggedright\arraybackslash}p{3.1cm} >{\raggedright\arraybackslash}p{3.0cm} L L@{}}
\toprule
\textbf{Bộ dữ liệu} & \textbf{Ngôn ngữ · miền} & \textbf{Quy mô} & \textbf{Nhiệm vụ} \\
\midrule
VLQA / DRILL~\cite{vlqa,drill} & Tiếng Việt · 27 lĩnh vực & 3.129 câu hỏi · 59.636 điều luật của 2.162 văn bản & Truy hồi điều luật và sinh câu trả lời có trích dẫn \\
MLQA-TSR, VLSP 2025~\cite{mlqatsr} & Tiếng Việt · biển báo GTĐB & 776 câu hỏi · 402 điều · 498 ảnh biển báo & Truy hồi và hỏi đáp đa phương thức ảnh -- văn bản \\
ALQAC~\cite{alqac} & Tiếng Việt (và tiếng Thái) & Chưa kiểm chứng từ nguồn chính thức & Truy hồi văn bản luật, suy luận kéo theo, hỏi đáp \\
Zalo AI Challenge 2021~\cite{zalo} & Tiếng Việt · văn bản luật & Chưa kiểm chứng từ nguồn chính thức & Truy hồi văn bản luật (Acc@1) \\
COLIEE~\cite{coliee} & Tiếng Nhật / Anh · dân sự, án lệ & 768 điều Bộ luật Dân sự Nhật (Task 3--4) & Truy hồi luật thành văn và kiểm tra kéo theo \\
\midrule
\textbf{Bộ dữ liệu của nhóm} & Tiếng Việt · GTĐB 2024--2026 & 120 câu hỏi · 345 hành vi · 109 quy tắc · 73 khái niệm & Tra cứu tri thức 7 lớp, gán nhãn tới \emph{định danh} tri thức \\
\bottomrule
\end{tabularx}
\end{table}

Bảng~\ref{tab:datasets} chỉ ghi con số khi kiểm chứng được từ bài báo hoặc trang chính thức;
các ô ``chưa kiểm chứng'' được để nguyên thay vì điền số ước đoán. Điểm cần lưu ý là
\textbf{đơn vị gán nhãn}: các bộ dữ liệu hiện có gán nhãn ở mức điều luật, trong khi bộ dữ
liệu của nhóm gán tới định danh của từng mẩu tri thức (một khái niệm, một quy tắc, một hành
vi kèm chế tài) --- nhờ vậy mới đo được Top-1 ở mức tri thức chứ không chỉ mức văn bản.

\subsection{Khoảng trống và chỗ đứng của đồ án}

Đối chiếu bốn hướng trên với yêu cầu Đề tài 4, nhóm xác định năm khoảng trống --- và đây
chính là chỗ đồ án đứng vào:

\begin{enumerate}
  \item \textbf{Chưa có hệ tra cứu tri thức cho GTĐB Việt Nam theo khung pháp lý 2024--2026.}
        Công trình đồ thị tri thức cùng miền gần nhất~\cite{dangKG} dựa trên Luật 23/2008 đã
        hết hiệu lực; MLQA-TSR~\cite{mlqatsr} tuy đúng khung luật mới nhưng giải bài toán
        nhận dạng biển báo đa phương thức, không phải tra cứu tri thức pháp lý.
  \item \textbf{Đầu ra là đoạn văn bản, không phải tri thức có cấu trúc.}
        VLQA/DRILL~\cite{vlqa,drill}, Zalo~\cite{zalo} và COLIEE Task~3~\cite{coliee} đều
        chấm ở mức ``xếp hạng điều luật đúng''; không nhiệm vụ nào yêu cầu trả về khái niệm,
        quan hệ và quy tắc suy diễn --- nên không đo được khả năng giải trình.
  \item \textbf{Chưa có công trình tiếng Việt kết hợp mô hình tri thức kiểu $K$ với suy diễn
        số học theo luật.} Legal-Onto~\cite{legalonto} và các mở
        rộng~\cite{phamLabor,phamLegalOntoUpdate,phamUnified} dừng ở tra cứu thuật ngữ và thủ
        tục, trong khi GTĐB đòi hỏi tính ra mức phạt tiền, thời hạn tước giấy phép và điểm
        trừ theo điều kiện phương tiện -- hành vi -- tình tiết.
  \item \textbf{Chưa mô hình hoá hiệu lực theo thời gian cho pháp luật Việt Nam.} Lý thuyết
        đã có~\cite{governatori,palmirani,legalruleml,demartim} nhưng phía Việt Nam mới dừng ở
        phát hiện thay đổi giữa hai phiên bản~\cite{phamLegalOntoUpdate}. Miền GTĐB thay đổi
        rất nhanh nên đây là khoảng trống có giá trị thực tiễn cao.
  \item \textbf{Hướng RAG/LLM chưa bảo đảm trung thực căn cứ} ở mức chấp nhận được cho tra
        cứu pháp luật~\cite{dahl,magesh,alce}. Đồ án chọn loại bỏ hoàn toàn khâu sinh, đổi độ
        ``mượt'' của câu trả lời lấy khả năng truy vết 100\% căn cứ --- một đánh đổi có chủ
        đích chứ không phải hạn chế kỹ thuật.
\end{enumerate}
